\scanpage{6}
by $\rho$, but not of $\unsure{\mathfrak S_3|\mathcal S_3}$, since we have for example
\[
\inserted{\sigma_0\sigma_1\bigl(=\sigma_1\sigma_\infty=\sigma_\infty\sigma_0\bigr)=\rho}
\]
\[
\deleted{[\sigma_0,\sigma_1]\unsure{\sigma_0^{-1}\sigma_1^{-1}}=(\sigma_0\sigma_1)^2=\rho^2}\qquad\text{in }\unsure{\mathfrak S_3|\mathcal S_3}
\]
but in $\operatorname{Aut}(\pi)$, \deleted{$\sigma_0\sigma_1$}
\[
\sigma_0\sigma_1(l_0)=\sigma_0(l_\infty)=l_1=\rho(l_0)
\]
\margin{$\unsure{2|z|t|\tau}$}
\[
\begin{aligned}
\sigma_0\sigma_1(l_1)&=\sigma_0(l_\infty^{-1}l_0^{-1})=\sigma_0(l_\infty)^{-1}\sigma_0\unsure{(}l_0)^{-1}\\
&=l_1^{-1}(l_1^{-1}l_\infty^{-1})^{-1}=l_1^{-1}l_\infty l_1
\end{aligned}
\]
\[
\begin{aligned}
\sigma_0\sigma_1(l_\infty)&=(l_1^{-1}l_\infty l_1)^{-1}l_1^{-1}=l_1^{-1}l_\infty^{-1}\deleted{\unsure{l_1|l_1^2|l_0^2}}\\
&=l_1^{-1}(l_\infty^{-1}l_1^{-1})l_1=l_1^{-1}l_0l_1
\end{aligned}
\]
thus $\sigma_0\sigma_1=\operatorname{int}(l_1^{-1})\rho$ in $\operatorname{Aut}(\pi)$
and likewise
\[
\sigma_1\sigma_\infty=\operatorname{int}(l_\infty^{-1})\rho
\]
\[
\sigma_\infty\sigma_0=\operatorname{int}(l_0^{-1})\rho
\]

We have a representation of $\unsure{\mathfrak S_3|\mathcal S_3}$ in the \uline{outer} group with loops only. Note that we cannot hope to realize this outer action by a genuine action (at least \unsure{not on|for|by on} the discrete group) since $\unsure{\mathfrak S_3|\mathcal S_3}$ acting on $\mathbb P^1_{\mathbb C}\setminus\{0,1,\infty\}$ has no fixed point. The fact that we have been able to make $\mathbb Z/3\mathbb Z$ act \inserted{on $\unsure{\pi|\Pi}$ itself} \unsure{comes|comes} from the fact that its action on $\mathbb P^1_{\mathbb C}-\{0,1,\infty\}$ admits one (and even two!) fixed point\inserted{, namely $\unsure{\pm\bar j|-\bar j|\pm j}$ ($j$ a primitive third root of $1$)} which we \deleted{\unsure{will be able to}} \unsure{have} taken as base point to compute the $\pi_1$

Let us also point out that we shall write in the Teichmüller group \inserted{$\Gamma_{0,3}$ acting on $\Pi_{0,3}$} the ``complex conjugation''
\mspage{5}
